\section{Preliminaries}
\label{sec:preliminaries}

In this work, we study the problem of generating differentially private explanations
for semantic differences between query outputs.

This section introduces the formal concepts and notation used throughout the paper.

\subsection{Databases and Queries}
\label{subsec:prelim-db-queries}

Let $D$ be a relational database over schema $\mathcal{R}$~\cite{Foundations}.
Each tuple $t \in D$ corresponds to an individual entity or record.
We denote by $Attr(D)$ the set of attributes appearing in the schema of $D$.

We consider SQL-style selection-projection queries and aggregate queries
with Boolean selection predicates over a single relational table
$D$ with schema $\mathcal{A}=(A_1,\ldots,A_m)$. Each attribute $A_i$
has an active domain $Dom(A_i)$, which may be categorical or numerical,
and each tuple $t \in D$ assigns a value $t[A_i] \in Dom(A_i)$ to every
attribute $A_i$. A query $Q$ consists of a selection predicate $\varphi$
over the attributes of $D$ and either returns the selected tuples or
computes an aggregate over them. Aggregate queries have the form
\[
Q(D) =
\texttt{SELECT } agg(A)
\texttt{ FROM } D
\texttt{ WHERE } \varphi,
\]
where $agg \in \{\texttt{COUNT},\texttt{SUM},\texttt{AVG}\}$ and
$\varphi$ is a Boolean combination of atomic conditions. Atomic
conditions include numerical comparisons, such as $A \leq c$ or
$A \geq c$, categorical equality predicates, such as $A = v$, and
categorical set-membership predicates, such as $A \in S$. Given a query
$Q$ and a database $D$, we denote by $Q(D)$ the output of the query
evaluated over $D$.



Throughout the paper, we study pairs of related queries:
\[
(Q_1,Q_2),
\]
where both queries are evaluated over the same database $D$.
Our goal is to explain the semantic differences between:
\[
Q_1(D)
\quad\text{and}\quad
Q_2(D).
\]

Throughout the paper, we study pairs of related query outputs, denoted
by $Q_1(D)$ and $Q_2(D)$, where both are derived from the same underlying
database $D$. The next subsection describes two application settings that
naturally give rise to such pairs.

\subsection{Application Settings: Query Refinement and Consistent Query Answering}

Query refinement studies the problem of modifying a user query so that
its output satisfies a desired set of constraints while remaining close
to the intent of the original query. In diversity-aware and fairness-aware
query processing, these constraints often specify required cardinalities
or proportions for groups defined by demographic or sensitive attributes.
For example, a refinement may relax, tighten, remove, or replace query
predicates in order to ensure that the result contains sufficient
representation of a given group \cite{MK09,LMSJ23,ERICA,CSSM24}.
Formally, given a database $D$, a query $Q$, and a set of output
constraints $\mathcal{C}$, a query-refinement mechanism returns a refined
query $Q'$ such that $Q'(D)$ satisfies $\mathcal{C}$, typically while
minimizing some notion of deviation from $Q$.

Consistent query answering addresses query evaluation over databases
that may violate a set of integrity constraints. Instead of evaluating a
query only over the possibly inconsistent database, CQA considers
repairs: database instances that satisfy the integrity constraints and
differ minimally from the original database. An answer is considered
consistent if it is obtained under the repair semantics, e.g., if it is
returned across the relevant repairs \cite{arenas1999,bertossi2006,bertossi2011}.
As a result, the answer produced by CQA may differ from the answer
obtained by evaluating the same query directly on the original database.

Both settings naturally induce the type of related query outputs studied
in this paper. In query refinement, the pair is the original and refined
query outputs, $(Q(D),Q'(D))$. In CQA, the pair consists of the direct
query answer over the original database and the answer obtained under
the repair semantics. Our goal is not to compute refinements, repairs, or
consistent answers, but to explain the semantic differences between such
related outputs in a privacy-preserving manner.

\subsection{Predicate Views}
\label{subsec:prelim-views}

To explain semantic differences between query outputs, our framework analyzes subsets
of the data defined by predicates. Predicate-defined subpopulations are commonly used
in explanation and query-analysis frameworks to localize meaningful changes in the
data~\cite{FEDEX,dpxplain2022,wu2013}.

\begin{definition}[Predicate View]
Given a database $D$ and a predicate $p$, the predicate view induced by $p$ is defined as:
\[
p(D)=\{t\in D \mid p(t)=\texttt{true}\}.
\]
\end{definition}

Predicate views allow the framework to focus on specific semantic regions of the data.

\subsection{Histograms and Marginals}
\label{subsec:prelim-histograms}

Our explanations are based on distributions of attributes over query outputs and predicate views.

\begin{definition}[Histogram]
Given a database $D$ and an attribute $A \in Attr(D)$, the histogram of $A$ over $D$
is the function:
\[
H_{D,A}: Dom(A,D)\rightarrow \mathbb{N},
\]
where $H_{D,A}(v)$ denotes the number of tuples in $D$ whose value for attribute $A$
equals $v$.
\end{definition}

In the standard statistical sense, a marginal distribution is obtained by
projecting a joint distribution onto a subset of attributes. In this
paper, we use one-way marginals, namely marginals over a single
attribute. Concretely, for a query output $Q(D)$ and an attribute $A$,
the one-way marginal of $A$ is the histogram $H_{Q(D),A}$ over all
values, of $A$. These one-way marginals are the basic
objects used in our explanations: comparing $H_{Q_1(D),A}$ and
$H_{Q_2(D),A}$ summarizes how the two query outputs differ with respect to attribute $A$. Such distributional
summaries are also widely used in privacy-preserving data analysis and release
mechanisms~\cite{dwork2014}.

For ease of presentation, we assume that the attribute values used in
histograms and predicate views are categorical. Numerical attributes can
be incorporated by rendering them categorical, e.g., by bucketizing their
domains into ranges, a common practice when presenting aggregate
results. Thus, for categorical attributes the histogram bins are the
attribute values, while for numerical attributes the bins correspond to
the chosen ranges.

\ag{No need to have a formal def for this. Just explain in text and add motivation for this} \alon{Addressed: removed the formal Bucketing definition and replaced it with
an explanation and motivation for using buckets for numerical
attributes.}

\ag{Unify with previous section about statistics.}\alon{Addressed}

To quantify semantic differences between query outputs, we use statistical distance
measures between distributions.

\begin{definition}[Total Variation Distance]
Given two distributions $P$ and $Q$ over the same domain $\Omega$, the Total Variation
Distance (TVD) between them is defined as:
\[
TVD(P,Q)
=
\frac{1}{2}
\sum_{x\in\Omega}
|P(x)-Q(x)|.
\]
\end{definition}

TVD is commonly used to measure similarity between distributions and is used throughout
our framework for view filtering and representativeness analysis.

\subsection{Differential Privacy}
\label{subsec:prelim-dp}
Because explanations are derived from sensitive data, our framework
relies on differential privacy to ensure that released explanations do
not reveal information about any particular individual~\cite{dwork2014}. Intuitively, a
differentially private mechanism should behave almost the same on two
databases that differ in one tuple. Thus, whether a single individual is
included in the database should have only a limited effect on the
distribution of the released output.

We first define neighboring datasets, which formalize the notion of two
databases differing in one individual. We then define differential
privacy as a bound on how much the output distribution of a randomized
mechanism can change between neighboring datasets. The parameter
$\epsilon$ controls the privacy loss: smaller values of $\epsilon$
provide stronger privacy but typically require adding more noise.

The amount of noise needed depends on the sensitivity of the function
being released, namely the maximum amount by which the function value can
change between neighboring datasets. The Laplace mechanism releases
noisy numerical statistics by adding noise calibrated to this sensitivity.
For private selection tasks, such as choosing informative views or
explanation candidates, we use the exponential mechanism, which selects
high-utility candidates with probability increasing in their utility
while still satisfying differential privacy.

\begin{definition}[Neighboring Datasets~\cite{dwork2006,dwork2014}]
Two datasets $D,D'$ are neighboring datasets if they differ by exactly one tuple.
We denote this relation by:
\[
D\sim D'.
\]
\end{definition}

\begin{definition}[Differential Privacy~\cite{dwork2006,dwork2014}]
A randomized mechanism $\mathcal{M}$ satisfies $\epsilon$-differential privacy if for
every pair of neighboring datasets $D\sim D'$ and every subset of outputs $S$:
\[
Pr[\mathcal{M}(D)\in S]
\leq
\exp(\epsilon)
Pr[\mathcal{M}(D')\in S].
\]
\end{definition}

\begin{definition}[Global Sensitivity~\cite{dwork2014}]
The global sensitivity of a function $f:D\rightarrow\mathbb{R}$ is defined as:
\[
\Delta f
=
\max_{D\sim D'}
|f(D)-f(D')|.
\]
\end{definition}

\begin{definition}[Laplace Mechanism~\cite{dwork2006,dwork2014}]
Given a function $f$ with global sensitivity $\Delta f$, the Laplace mechanism releases:
\[
f(D)+Laplace(\Delta f/\epsilon).
\]
\end{definition}

Thus, functions with lower sensitivity require less noise for the same
privacy budget, which motivates the low-sensitivity proxy scores used in
our algorithm.

\begin{definition}[Exponential Mechanism~\cite{mcsherry2007,dwork2014}]
Given a dataset $D$, a set of candidate outputs $Y$, a utility function
$u:Y\times D\rightarrow\mathbb{R}$, and privacy budget $\epsilon$, the exponential
mechanism selects an output $y\in Y$ with probability proportional to:
\[
\exp\left(
\frac{\epsilon u(y,D)}
{2\Delta u}
\right),
\]
where $\Delta u$ is the global sensitivity of the utility function.
\end{definition}

The exponential mechanism enables private selection of high-quality explanation
candidates while preserving differential privacy.


Differential privacy satisfies several important composition properties, enabling
complex privacy-preserving analytical workflows constructed from multiple private
mechanisms \cite{mcsherry2009,li2016,bun2016}.
These properties are central to the privacy analysis of our algorithms, which combine
private view selection, private candidate selection, and noisy histogram release.

\begin{theorem}[Sequential Composition~\cite{dwork2014,li2016}]
Let $\mathcal{M}_1,\ldots,\mathcal{M}_k$ be randomized mechanisms such that
$\mathcal{M}_i$ satisfies $\epsilon_i$-differential privacy for every
$i \in \{1,\ldots,k\}$.
Then the mechanism that releases the joint output
\[
\mathcal{M}(D)
=
\left(
\mathcal{M}_1(D),
\ldots,
\mathcal{M}_k(D)
\right)
\]
satisfies
\[
\left(\sum_{i=1}^{k}\epsilon_i\right)
\text{-differential privacy}.
\]
\end{theorem}

Sequential composition allows us to account for the total privacy loss of an algorithm
that invokes multiple differentially private subroutines by summing their individual
privacy budgets.

\begin{theorem}[Post-Processing~\cite{dwork2014,li2016}]
Let $\mathcal{M}$ be an $\epsilon$-differentially private mechanism,
and let $f$ be any randomized mapping that does not access the private database
except through the output of $\mathcal{M}$.
Then the mechanism
\[
f(\mathcal{M}(D))
\]
also satisfies $\epsilon$-differential privacy.
\end{theorem}

The post-processing property implies that any deterministic or randomized computation
performed only on already-private outputs does not incur additional privacy loss.

Differential privacy is also preserved under adaptive composition
\cite{mcsherry2009,li2016}.
That is, if each mechanism in a sequence is selected as a function
of previously released outputs, then the overall privacy loss remains
bounded by the sum of the individual privacy budgets.
